This dissertation argues that the tools through which we create and perceive sound carry within them specific cultural orientations. Digital audio systems, while presented as universal, encode the temperaments and assumptions of their origins. Equal temperament tuning, grid-based sequencing, and Eurocentric standards of fidelity organize perception according to a narrow acoustic worldview. These defaults do more than shape what can be heard; they shape what can be imagined.

When Ruha Benjamin warns that "the road to inequity is paved with technical fixes" \parencite[p. 13]{Benjamin_2019a}, she identifies why surface-level diversity initiatives often fail to address the foundational issue: technical fixes assume problems lie in biased implementations of otherwise neutral procedures, maintaining what Sylvia Wynter calls the over-representation of "Man," the Western Bourgeois conception that has "positioned itself as equivalent to humanity itself" through supposedly universal logics \parencite{Wynter_2003}. 

The framework of \textit{Speculocultural Technopoiesis} responds to this condition by proposing a form of sonic research-creation grounded in modification. Through iterative cycles of making, listening, and reflection, we explore how culturally informed approaches to technology can generate new sonic and epistemic possibilities. In this view, modification functions as a form of \textit{transduction}: speculative cultural ideas are converted into material form through the redesign of software, instruments, and interfaces. The resulting shifts in perception, what we can describe as \textit{refraction}, make audible relationships that were previously muted by the defaults of dominant design.

By engaging in this way,we can position Black practices (around sound and material) as catalysts for particular insights within this form of  computational navigation; precisely because they have historically operated through tactical appropriation and transformation of dominant systems. Black (and other Afro-diasporic cultures) creative practices have consistently demonstrated how the same tools produce fundamentally different possibilities when approached through diametric ways of knowing. These practices affirm (in their own way) that tools cannot be a neutral container. They function, fundamentally, as systems that morph the shape and function of their mediated perception when operated through distinct hands.

Within this recognition, \textit{Speculocultural Technopoiesis} situates sound as a site of both philosophical inquiry and practical intervention, where cultural imagination becomes a technical principle. By engaging with the non-neutrality and the non-universality of technologies identified by Ihde and Benjamin \parencite{Ihde_1990, Benjamin_2019a}, this research demonstrates how altering the structure of tools can also alter the structure of thought that stems from those same tools.

The central claim is that sonic modification can operate as a mode of epistemic resistance: a way of revealing and re-tuning the cultural frequencies embedded within our technologies. Through this process, \textit{Speculocultural Technopoiesis} seeks to develop new approaches to digital audio design that emerge from Black and Afro-diasporic sonic epistemologies. Those are the approaches that imagine and materialize those "other" relationships between sound, culture, and computation. The following section outlines the specific research questions that are guiding the basis of this inquiry.